\documentclass[11pt]{report}

% Shared preamble for main.tex and Proof_Attempts.tex
% (packages, theorem styles, macros, page layout, etc.)

% -------------------------------------------------
% Page layout / typography
% -------------------------------------------------
\usepackage[a4paper,margin=30mm,headheight=14pt]{geometry}
\usepackage{microtype}              % better spacing/kerning
\usepackage[T1]{fontenc}
\usepackage[utf8]{inputenc}         % ok for pdflatex; remove if using lualatex/xelatex
\usepackage{xcolor}
\usepackage{lmodern}                % modern Latin Modern fonts
\usepackage{setspace}               % optional: \onehalfspacing etc.

% -------------------------------------------------
% Graphics / tables / lists
% -------------------------------------------------
\usepackage{graphicx}
\usepackage{booktabs}               % professional tables
\usepackage{array}
\usepackage{multirow}
\usepackage{caption}
\usepackage{subcaption}
\usepackage{enumitem}

% -------------------------------------------------
% Math (standard praxis)
% -------------------------------------------------
\usepackage{amsmath,amssymb,amsthm} % AMS core
\usepackage{mathtools}             % extends/fixes amsmath
\usepackage{stmaryrd}              % \llbracket, \rrbracket for stochastic intervals

% Upright differential (use as \int f(x)\,\diff x)
\newcommand{\diff}{\mathop{}\!\mathrm{d}}

% Operators
\DeclareMathOperator*{\arginf}{arginf}
\DeclareMathOperator*{\argmin}{arg\,min}

\usepackage{bm}                    % bold math symbols
\usepackage{dsfont}                % \mathds{1} indicator
\usepackage{mathrsfs}              % script font \mathscr
\usepackage{bbm}                   % alternative blackboard bold, incl. \mathbbm{1}
\usepackage{siunitx}               % units/numbers (also nice in finance tables)

% -------------------------------------------------
% References / hyperlinks (standard modern setup)
% -------------------------------------------------
\usepackage{xcolor}
\usepackage{hyperref}
\hypersetup{
  colorlinks=true,
  linkcolor=blue,
  citecolor=blue,
  urlcolor=blue,
  pdftitle={Master Project},
  pdfauthor={Your Name},
  pdfsubject={Mathematical Finance},
  pdfkeywords={Mathematical Finance, Stochastic Calculus, ...}
}

% cleveref + shared theorem counters:
% If environments share the "theorem" counter, \cref may otherwise call everything "Theorem".
% aliascnt creates alias counters so cleveref can distinguish Remark/Lemma/etc.
\usepackage{aliascnt}

% cleveref must be loaded AFTER hyperref
\usepackage[nameinlink,noabbrev]{cleveref}

% -------------------------------------------------
% Theorems / definitions (standard mathematical setup)
% -------------------------------------------------
% Numbering: Theorem 2.3, Lemma 2.4, ... within each chapter.
\theoremstyle{plain} % italic body, bold heading
\newtheorem{theorem}{Theorem}[chapter]

\newaliascnt{proposition}{theorem}
\newtheorem{proposition}[proposition]{Proposition}
\aliascntresetthe{proposition}

\newaliascnt{lemma}{theorem}
\newtheorem{lemma}[lemma]{Lemma}
\aliascntresetthe{lemma}

\newaliascnt{corollary}{theorem}
\newtheorem{corollary}[corollary]{Corollary}
\aliascntresetthe{corollary}

\theoremstyle{definition} % upright body
\newaliascnt{definition}{theorem}
\newtheorem{definition}[definition]{Definition}
\aliascntresetthe{definition}

\newaliascnt{assumption}{theorem}
\newtheorem{assumption}[assumption]{Assumption}
\aliascntresetthe{assumption}

\theoremstyle{remark} % upright body, italic heading
\newaliascnt{remark}{theorem}
\newtheorem{remark}[remark]{Remark}
\aliascntresetthe{remark}

% Unnumbered variants (use sparingly)
\newtheorem*{theorem*}{Theorem}
\newtheorem*{lemma*}{Lemma}
\newtheorem*{definition*}{Definition}

% Cleveref names for theorem-like environments
\crefname{theorem}{Theorem}{Theorems}
\Crefname{theorem}{Theorem}{Theorems}
\crefname{proposition}{Proposition}{Propositions}
\Crefname{proposition}{Proposition}{Propositions}
\crefname{lemma}{Lemma}{Lemmas}
\Crefname{lemma}{Lemma}{Lemmas}
\crefname{corollary}{Corollary}{Corollaries}
\Crefname{corollary}{Corollary}{Corollaries}
\crefname{definition}{Definition}{Definitions}
\Crefname{definition}{Definition}{Definitions}
\crefname{assumption}{Assumption}{Assumptions}
\Crefname{assumption}{Assumption}{Assumptions}
\crefname{remark}{Remark}{Remarks}
\Crefname{remark}{Remark}{Remarks}

% -------------------------------------------------
% Bibliography (biblatex + biber)
% -------------------------------------------------
\usepackage{csquotes} % recommended with biblatex
\usepackage[
  backend=biber,
  style=authoryear,
  dashed=false, % don't replace repeated author/editor with an em-dash
  maxcitenames=2,
  maxbibnames=99
]{biblatex}
\addbibresource{references.bib}

% -------------------------------------------------
% Appendix + PDF inclusion (optional)
% -------------------------------------------------
\usepackage{appendix}
\usepackage{pdfpages}

% -------------------------------------------------
% Nomenclature / symbols list (optional but common in math finance)
% -------------------------------------------------
\usepackage[intoc]{nomencl}
\makenomenclature
% Example: \nomenclature{$W_t$}{Standard Brownian motion}

% -------------------------------------------------
% Header / footer (simple, thesis-friendly)
% -------------------------------------------------
\usepackage{fancyhdr}
\pagestyle{fancy}
\fancyhf{}
\fancyhead[LE,RO]{\thepage}
\fancyhead[LO]{\nouppercase{\rightmark}}
\fancyhead[RE]{\nouppercase{\leftmark}}
\renewcommand{\headrulewidth}{0.4pt}

% -------------------------------------------------
% Table of contents depth (DETAILED ToC)
% -------------------------------------------------
% tocdepth: 0=part, 1=chapter, 2=section, 3=subsection, 4=subsubsection
\setcounter{tocdepth}{3}  % show down to \subsection in ToC
\setcounter{secnumdepth}{3} % number down to \subsection (or 4 for subsubsection)

\begin{document}
\section{Completentess of \(L^2(A)\cap L^2(M)\)}
Let \(X=M+A\) be the canonical decomposition of \(X\in\mathcal
S^2_{loc}\). Let 
\begin{equation}
    L^2(M) = \left\{\vartheta\text{ predictable }:\mathbb{E}[\int_0^T\vartheta_s^2\diff [M]_s]=\lVert\vartheta\rVert_{L^2(M)}^2<\infty \right\}
\end{equation}
 and 
 \begin{equation}
     L^2(A) = \left\{\vartheta\text{ predictable }:\mathbb{E}[\left(\int_0^T|\vartheta_s|\diff |A|_s\right)^2]=\lVert\vartheta\rVert_{L^2(A)}<\infty \right\}
 \end{equation}
 then 
 \begin{equation}
     \Theta=L^2(M)\cap L^2(A),\;\lVert\cdot\rVert_{\Theta}=\lVert\cdot\rVert_{L^2(M)}+\lVert\cdot\rVert_{L^2(A)}
 \end{equation}
 In order to show that \(\Theta\) is complete under the norm \(\lVert\cdot\rVert_{\Theta}\), it is enough that \((L^2(M),\lVert\cdot\rVert_{L^2(M)})\) and \((L^2(A),\lVert\cdot\rVert_{L^2(A)})\) are complete normed spaces. 
 \\\\
 Q1: Why? If \((\vartheta_n)_{n\in\mathbf{N}}\) is a Cauchy sequence in \((\Theta,\lVert\cdot\rVert_{\Theta})\) then it is a Cauchy sequence in \((L^2(M),\lVert\cdot\rVert_{L^2(M)})\) and \((L^2(A),\lVert\cdot\rVert_{L^2(A)})\). If the latter have been shown to be complete under their norm, there exists \(\vartheta^*\) such that \(\vartheta_n\overset{L^2(M)}{\to}\vartheta^*\) and \(\widetilde{\vartheta}^*\) such that \(\vartheta_n\overset{L^2(A)}{\to}\widetilde{\vartheta}^*\). The claim is shown if \(\lVert\vartheta^*-\widetilde{\vartheta}^*\rVert_{\Theta}=0\). I don't see why this should be the case.
 \\\\
 Assume the latter to be true. The completeness of \(L^2(M)\) is clear as it is isomorphic to \(L^2(\mathcal{P},\mathbb{P}_M)\) and the latter is complete (Jean-Francois Le Gall: Measure Theory, Probability and Stochastic Processes, Thm.4.5). But for the case \(L^2(A)\) it is not clear to me. In particular, the standard \(L^p\) completeness arguments do not work because \(L^2(A)\) is not isomorphic to an \(L^p\)-space. Indeed, the square in the definition of \(\lVert\cdot\rVert_{L^2(A)}\) destroys subadditivity, so that \(B\mapsto \mathbb{E}[\left(\int_0^T1_B\diff|A|\right)^2]\) does not define a measure on \((\Omega\times[0,T],\mathcal{P})\). (In the case \(p=1\) it does and the standard argument goes through).
 \bigbreak
 \underline{Attempt 1:}
 \\
 Assume it is shown that \((L^1(A),\lVert\cdot\rVert_{L^1(A)})\) is complete. 
 \\
 Assume \(\sum_{n\geq 1}\lVert \vartheta^n\rVert_{L^2(A)}<\infty\). Note that this equivalently reads
 \begin{equation}
     \sum_{n\geq 1}\lVert\int_0^T|\vartheta^n|\diff |A|\rVert_{L^2(P)}<\infty
 \end{equation}
 We denote \(\theta^n = \sum_{i=1}^n\vartheta^i\).
 \\
 Because \(L^2(P)\) is complete, there exists \(Z_2\) such that 
 \begin{equation}
     \lim_{n\to\infty}\lVert\int_0^T|\theta^n|\diff |A|-Z_2\rVert_{L^2(P)}=0\;.
 \end{equation}
 Furthermore, as
 \begin{equation}
     \lVert f\rVert_{L^1(P)}\leq P(\Omega)^{\frac{1}{2}}\lVert f\rVert_{L^2(P)}
 \end{equation}
 \begin{equation}
     \sum_{n\geq 1}\lVert \vartheta^n\rVert_{L^2(A)}<\infty\implies\sum_{n\geq 1}\lVert \vartheta^n\rVert_{L^1(A)}<\infty
 \end{equation}
 Again, \(L^1(P)\) is well known to be complete, hence \(\exists Z_1\) such that
  \begin{equation}
     \lim_{n\to\infty}\lVert\int_0^T|\theta^n|\diff |A|-Z_1\rVert_{L^1(P)}=0\;.
 \end{equation}
 Furthermore, we assumed that \((L^1(A),\lVert\cdot\rVert_{L^1(A)})\) is known to be complete. Hence, there exists \(\theta^*\) such that 
 \begin{equation}
     \lim_{n\to\infty}\mathbb{E}[\int_0^\infty|\theta^n-\theta^*|\diff |A|]=0
 \end{equation}
 Because \(|\theta^n-\theta^*|\geq ||\theta^n|-|\theta^*||\) the latter implies
 \begin{equation}
     \lim_{n\to\infty}\mathbb{E}[\left|\int_0^T|\theta^n|\diff |A|-\int_0^T|\theta^*|\diff |A|\right|]=0
 \end{equation}
 that is \(\int_0^T|\theta^n|\diff |A|\to\int_0^T|\theta^*|\diff |A|\) in \(L^1\). \(L^1\), being a normed vector space, is in particular Hausdorff, hence has unique limits. Thus, 
 \begin{equation}
     \int_0^T|\theta^*|\diff |A|=Z_1
 \end{equation}
 But 
 \begin{equation}
     \lim_{n\to\infty}\lVert Z_2-\int_0^T|\theta^n|\diff |A|\rVert_{L^1}\lesssim \lVert Z_2-\int_0^T|\theta^n|\diff |A|\rVert_{L^2}=0
 \end{equation}
 Hence, again by the uniqueness of limits, 
 \begin{equation}
     Z_2=Z_1=\int_0^T|\theta^*|\diff |A|
 \end{equation}
 where the equalities hold in \(L^1\). 
 This shows that
 \begin{equation}
     \mathbb{E}[\left(\int_0^T|\theta^n|-|\theta^*|\diff |A|\right)^2]\to0
 \end{equation}
 \bigbreak
 \underline{Attempt 2:}
 \end{document}
\end{document}